% 第九卷 7.7.1 节：DeepSeek-R1-Zero与DeepSeek-R1的训练路径（跨栏，约160 mm宽）。
% 上排R1-Zero；下两行R1四阶段与蒸馏分支。
% 黄色：训练阶段；蓝色：模型或数据；红色：交付模型。实线箭头为参数继承，虚线箭头为数据迁移。
\begin{tikzpicture}[
  x=1mm,y=1mm,
  figure picture,
  model/.style={figure node,draw=FigureInput,fill=FigureInputFill,minimum height=9mm,
    minimum width=24mm,inner xsep=2pt,inner ysep=2pt},
  stage/.style={figure node,draw=FigureModel,fill=FigureModelFill,minimum height=9mm,
    minimum width=30mm,inner xsep=2pt,inner ysep=2pt,line width=.9pt},
  data/.style={model,fill=white},
  final/.style={figure node,draw=FigureOutput,fill=FigureOutputFill,minimum height=9mm,
    minimum width=16mm,inner xsep=2pt,inner ysep=2pt},
  move/.style={figure flow,dash pattern=on 2.5pt off 1.5pt,draw=FigureMuted}
]
  % ---------- R1-Zero ----------
  \draw[figure subgroup] (0,54) rectangle (160,76);
  \node[figure heading,anchor=north west] at (1.5,75) {DeepSeek-R1-Zero};
  \node[model] (b0) at (16,62) {DeepSeek-V3-Base};
  \node[stage] (z1) at (58,62) {GRPO强化学习\\正确性奖励与格式奖励};
  \node[final] (z2) at (100,62) {R1-Zero};
  \draw[figure flow] (b0) -- (z1);
  \draw[figure flow] (z1) -- (z2);
  \node[figure note,anchor=west,align=left] at (112,62) {模板：推理写在think标记内，\\答案写在answer标记内};

  % ---------- R1 ----------
  \draw[figure subgroup] (0,-12) rectangle (160,50);
  \node[figure heading,anchor=north west] at (1.5,49) {DeepSeek-R1};
  \node[model] (b1) at (16,38) {DeepSeek-V3-Base};
  \node[stage] (s1) at (56,38) {阶段1：冷启动SFT\\数千条长思维链};
  \node[stage] (s2) at (98,38) {阶段2：推理强化学习\\加语言一致性奖励};
  \node[data] (d) at (140,38) {阶段3：拒绝采样\\约60万推理+20万通用};
  \draw[figure flow] (b1) -- (s1);
  \draw[figure flow] (s1) -- (s2);
  \draw[move] (s2) -- (d);

  \node[model] (b2) at (16,14) {DeepSeek-V3-Base};
  \node[stage] (s3) at (56,14) {阶段3：从基座再SFT\\两类数据合并};
  \node[stage] (s4) at (98,14) {阶段4：全场景强化学习\\规则奖励与偏好奖励};
  \node[final] (r1) at (134,14) {R1};
  \draw[figure flow] (b2) -- (s3);
  \draw[figure flow] (s3) -- (s4);
  \draw[figure flow] (s4) -- (r1);
  \draw[move] (d.south) -- ++(0,-6) -| (s3.north);

  \node[stage,minimum width=34mm] (kd) at (126,-4) {蒸馏：较小模型SFT\\Qwen、Llama学生};
  \draw[move] ([xshift=12mm]d.south) |- (kd.east);
\end{tikzpicture}
